% 112-75(112쪽) — 삼차함수 y=f(x) 의 그래프(x축과 x=-1, 1, 3 에서 만나고 y절편 -3 · 최고차항 계수 음수). 스캔 화질이 낮아 TikZ 로 다시 그림.
% 원문에 있는 것만: 눈금 숫자 -1 · 1 · 3 (x) · -3 (y절편) · 라벨 O, x, y, y=f(x). 극값의 좌표는 원문에 없어 적지 않는다.
\begin{tikzpicture}[x=0.42cm, y=0.26cm, line width=0.5pt]
  \draw[->, >=stealth, line width=0.7pt] (-1.9,0) -- (4.2,0) node[below=1pt] {$x$};
  \draw[->, >=stealth, line width=0.7pt] (0,-5.8) -- (0,4.8) node[left=1pt] {$y$};
  \node[below left=-1pt and -2pt, font=\small] at (0,0) {$\pt{O}$};
  \draw[line width=0.6pt, domain=-1.35:3.5, samples=200, smooth] plot (\x, {-(\x+1)*(\x-1)*(\x-3)});
  \node[below=3pt, font=\small] at (-1.32,0) {$-1$};
  \node[below=3pt, font=\small] at (1.25,0) {$1$};
  \node[below=3pt, font=\small] at (2.8,0) {$3$};
  \node[right=2pt, font=\small] at (0,-3) {$-3$};
  \node[right=1pt, font=\small] at (1.4,4.4) {$y=f(x)$};
\end{tikzpicture}
